\documentclass[11pt]{article}

\usepackage{amsmath}
\usepackage{amssymb}
\usepackage{amsthm}
\usepackage[margin=1in]{geometry}

\newtheorem{theorem}{Theorem}
\newtheorem{proposition}[theorem]{Proposition}
\theoremstyle{definition}
\newtheorem{definition}[theorem]{Definition}
\newtheorem{remark}[theorem]{Remark}

\title{A counterexample to conjecture 00000001676:\\ the isoperimetric profile of the torus grid}
\author{Submission \texttt{earthking11\_submission\_20260913120000}}
\date{2026-09-13}

\begin{document}
\maketitle

\begin{abstract}
Conjecture 00000001676 asserts that the isoperimetric profile of the
$n\times n$ torus grid $C_n \mathbin{\square} C_n$ is
$\operatorname{ip}(m)=\lceil 4\sqrt{nm-m^2}\rceil$ for all $m\le n^2/2$.
We refute the conjecture. It has two independent defects. First, on part of
its stated range the radicand $nm-m^2$ is negative, so the formula is not even
real: for $n=3$, $m=4$ one has $m\le n^2/2=4.5$ but $nm-m^2=-4$. Second,
wherever it is defined the formula is wrong. The cleanest counterexample is
$(n,m)=(3,3)$: the formula gives $\lceil 4\sqrt{9-9}\rceil=0$, whereas the true
minimum edge boundary is $6$.
\end{abstract}

\section{Definitions}

\begin{definition}[Torus grid]
Let $C_n$ be the cycle graph on $n$ vertices. The \emph{$n\times n$ torus grid}
is the Cartesian product $C_n \mathbin{\square} C_n$. Its vertex set is
$\mathbb{Z}_n\times\mathbb{Z}_n$, and two vertices are adjacent when they differ
in exactly one coordinate, differences being taken modulo $n$ (the torus
wrap-around). The graph is $4$-regular and has $n^2$ vertices and $2n^2$ edges.
\end{definition}

\begin{definition}[Edge boundary and isoperimetric profile]
For $S\subseteq V(C_n\mathbin{\square} C_n)$ the \emph{edge boundary} is
\[
  \partial S \;=\; \bigl\{\, \{u,v\}\in E : |\{u,v\}\cap S|=1 \,\bigr\},
  \qquad
  |\partial S|
  = \sum_{v\in S}\bigl(4-|N(v)\cap S|\bigr).
\]
The \emph{isoperimetric profile} is
$\operatorname{ip}(m)=\min\{\,|\partial S| : S\subseteq V,\ |S|=m\,\}$.
\end{definition}

\begin{definition}[Conjectured formula]\label{def:formula}
\emph{(Conjecture 00000001676.)} For all $m\le n^2/2$,
\[
  \operatorname{ip}(m) \;=\; \Bigl\lceil 4\sqrt{\,nm-m^{2}\,}\Bigr\rceil ,
\]
attained by diagonal cuts, with no other minimizing configuration.
\end{definition}

\section{Exhaustive computation of the true profile}

For $n=3$ and $n=4$ we computed $\operatorname{ip}(m)$ by brute force: enumerate
every $m$-vertex subset of the $n^2$ vertices, evaluate $|\partial S|$, and take
the minimum. Subsets were encoded as bitmasks and $\partial S$ was evaluated as
$\sum_{v\in S}\#\{w\in N(v): w\notin S\}$, which counts every crossing edge once.
For $C_3\mathbin{\square} C_3$ there are $2^9=512$ masks; for
$C_4\mathbin{\square} C_4$, $2^{16}=65536$ masks. The results are in
Table~\ref{tab:results}. (The script \texttt{reproduce.py} regenerates them; the
Lean development in \texttt{lean4/} verifies the decisive finite cases by kernel
evaluation over the $512$ masks.)

\begin{table}[htbp]
\centering
\begin{tabular}{c|c|c|c|l}
\multicolumn{5}{c}{\textbf{$C_3\mathbin{\square} C_3$, stated range $m\le 9/2$, i.e.\ $m=1,\dots,4$}}\\[2pt]
$m$ & $\operatorname{ip}(m)$ true & formula & $nm-m^2$ & status\\ \hline
$1$ & $4$ & $\lceil 4\sqrt2\rceil=6$ & $2$  & formula wrong\\
$2$ & $6$ & $6$                 & $2$  & agrees\\
$3$ & $6$ & $\lceil 4\sqrt0\rceil=0$ & $0$ & \textbf{formula wrong}\\
$4$ & $8$ & undefined           & $-4$ & \textbf{not real}\\
\end{tabular}

\vspace{1.2em}

\begin{tabular}{c|c|c|c|l}
\multicolumn{5}{c}{\textbf{$C_4\mathbin{\square} C_4$, stated range $m\le 8$}}\\[2pt]
$m$ & $\operatorname{ip}(m)$ true & formula & $nm-m^2$ & status\\ \hline
$1$ & $4$  & $\lceil 4\sqrt3\rceil=7$  & $3$   & formula wrong\\
$2$ & $6$  & $8$                      & $4$   & formula wrong\\
$3$ & $8$  & $7$                      & $3$   & formula wrong\\
$4$ & $8$  & $\lceil 4\sqrt0\rceil=0$  & $0$   & formula wrong\\
$5$ & $10$ & undefined                & $-5$  & not real\\
$6$ & $10$ & undefined                & $-12$ & not real\\
$7$ & $10$ & undefined                & $-21$ & not real\\
$8$ & $8$  & undefined                & $-32$ & not real\\
\end{tabular}
\caption{True minimum edge boundary versus the conjectured formula.}
\label{tab:results}
\end{table}

\section{The counterexample}

\begin{theorem}[Cleanest counterexample]\label{thm:main}
For the torus grid $C_3\mathbin{\square} C_3$ and $m=3$, the conjectured formula
of Definition~\ref{def:formula} gives $0$, but the true minimum edge boundary is
$6$. Hence the conjecture is false.
\end{theorem}

\begin{proof}
Here $n=m=3$, so $nm-m^2=9-9=0$ and the formula gives
$\lceil 4\sqrt{0}\rceil=0$.

A full row $\{(i,0),(i,1),(i,2)\}$ of the torus has $3$ vertices. Each of its
vertices has its two horizontal neighbours inside the row and its two vertical
neighbours outside the row, so the induced subgraph has $3$ edges and
\[
  |\partial S| = 3\cdot 4 - 2\cdot 3 = 12-6 = 6 .
\]
The exhaustive search over all $\binom{9}{3}=84$ triples (Table~\ref{tab:results})
shows no triple has a smaller boundary, so $\operatorname{ip}(3)=6$. Thus the
conjecture predicts $0$ where the truth is $6$.
\end{proof}

\begin{proposition}[Further defined-range failures]
The formula is also wrong at $(n,m)=(3,1)$: it gives
$\lceil 4\sqrt{2}\rceil=6$, while a single vertex has boundary $4$. More
generally Table~\ref{tab:results} exhibits six distinct defined values of
$(n,m)$ in the stated range at which the formula disagrees with the true
profile.
\end{proposition}

\section{Ill-posedness of the stated range}

The conjecture quantifies over all $m\le n^2/2$. But the radicand $nm-m^2
= m(n-m)$ is negative precisely when $m>n$. Together with the upper bound
$m\le n^2/2$ this is non-empty for every $n\ge 3$: choose any integer $m$ with
$n<m\le n^2/2$ (for $n=3$, $m=4$; for $n=4$, $m=5,6,7,8$; in general such an
$m$ exists because $n^2/2\ge n+1$ for $n\ge 3$). At those points the formula is
not a real number, so the conjecture is not even well posed on its stated
domain.

\begin{remark}
For $n=3$, $m=4$: $nm-m^2=12-16=-4<0$, while $m=4\le n^2/2=4.5$. For $n=4$,
$m=8$: $nm-m^2=32-64=-32<0$, while $m=8=n^2/2$. Thus the defect occurs within
the stated range, not merely beyond a boundary convention.
\end{remark}

\section{Verdict}

The conjecture is \textbf{false}, on two independent grounds:
\begin{enumerate}
  \item \textbf{Ill-posed on part of its range.} For $n<m\le n^2/2$ the
        radicand $nm-m^2$ is negative (e.g.\ $n=3$, $m=4$), so the formula is
        not real.
  \item \textbf{Wrong wherever defined.} The cleanest instance is
        $(n,m)=(3,3)$: the formula gives $0$ while the true minimum edge
        boundary is $6$. Other defined-range failures include $(n,m)=(3,1)$
        and $(4,1),(4,2),(4,3),(4,4)$.
\end{enumerate}

\end{document}
